% !TeX root = ../main.tex
% ---------------------------------------------------------------------------
% Water Inrush, Gas and Construction Hazards
% Scope extension: not derived from the RSE3010 weekly material.
% Key sources: BTS/ICE Tunnel Design and Construction guidance; AS 2865
% (confined spaces); Victorian tunnelling safety regulations; ITA WG5 on
% health and safety.
% ---------------------------------------------------------------------------
\chapter{Water Inrush, Gas and Construction Hazards}
\label{ch:construction-hazards}


\section{Hazard identification for underground works}
\label{sec:hazard-identification}
\scopenote{Systematic hazard identification along an alignment, and the
distinction between geotechnical hazards, construction hazards and
occupational hazards.}

\section{Water inrush and cavities}
\label{sec:water-inrush-and-cavities}

\subsection{Water inrush: mechanisms and precursors}
\label{sec:inrush-mechanisms}
\scopenote{Confined aquifers, karst conduits, fault-hosted flow paths and
buried valleys; observable precursors at the face and in probe holes.}

\subsection{Probing and pre-excavation treatment}
\label{sec:probing}
\scopenote{Probe drilling policy, standoff distance, controlled bleed-off, and
pre-excavation grouting as the standard response. Cross-reference to
Chapter~\ref{ch:ground-improvement}.}

\subsection{Cavity and karst detection}
\label{sec:karst-detection}
\scopenote{Surface and borehole geophysics, look-ahead seismic and
electromagnetic methods, and their practical detection limits ahead of a face.}

\section{Gas hazards}
\label{sec:gas-hazards-group}

\subsection{Gas hazards}
\label{sec:gas-hazards}
\scopenote{Methane from coal measures and organic soils, hydrogen sulfide,
carbon monoxide and dioxide, radon and oxygen deficiency; sources, migration
paths and accumulation points.}

\subsection{Gas classification, monitoring and controls}
\label{sec:gas-controls}
\scopenote{Classifying a tunnel as potentially gassy or gassy, fixed and
personal monitoring, intrinsically safe plant, ignition control, and the
regulatory regime in Australian practice.}

\section{Ventilation, heat and dust}
\label{sec:ventilation-heat-and-dust}

\subsection{Construction ventilation}
\label{sec:construction-ventilation}
\scopenote{Airflow requirements for dust, fumes, gas dilution and heat;
forcing and exhausting systems; duct losses and fan selection; ventilation as
a constraint on drive length.}

\subsection{Heat, humidity and thermal stress}
\label{sec:heat-stress}
\scopenote{Virgin rock temperature at depth, machine heat load, and cooling
requirements for worker safety in deep or long drives.}

\subsection{Dust and silica exposure}
\label{sec:dust}
\scopenote{Respirable crystalline silica, exposure standards, dust suppression
and extraction, and health surveillance obligations.}

\section{Contamination and fire}
\label{sec:contamination-and-fire}

\subsection{Contaminated and hazardous ground}
\label{sec:contaminated-ground}
\scopenote{Contaminated fill and groundwater, asbestos-bearing rock, acid
sulfate soils; handling, treatment and disposal of contaminated spoil.}

\subsection{Fire safety during construction}
\label{sec:construction-fire}
\scopenote{Ignition sources, fire load from plant and services, detection,
suppression, refuge chambers and self-rescuers.}

\section{Access, egress and emergency response}
\label{sec:access-egress-and-emergency-response}

\subsection{Confined space, access and egress}
\label{sec:confined-space}
\scopenote{Confined space classification, permit systems, secondary egress
requirements and the constraints of a single-entry drive.}

\subsection{Emergency planning and rescue}
\label{sec:emergency-planning}
\scopenote{Emergency response plans, tunnel rescue capability, communication
systems, and rehearsal. Lessons from documented incidents.}
